% HerHealthGPT-LU / LUHME 2026 — Introduction and Related Work
% Paste into ACL Overleaf template (double-blind safe).

\section{Introduction}

Women's-health communication is a high-stakes \emph{language understanding} problem.
Patients describe menstrual, hormonal, and fertility concerns through colloquial, indirect, emotionally loaded, and culturally coded language that is often clinically under-specified \cite{recker-etal-2025-large}.
A model that misreads such an utterance---confusing routine cycle variation with a fertility emergency, collapsing PCOS-related hormonal symptoms into unrelated categories, or reassuring when clarification is warranted---can steer a user toward delayed care, unnecessary anxiety, or unsafe self-management long before any downstream answer is generated.
Recent clinical commentary on large language models (LLMs) in gynecology and obstetrics therefore stresses not only whether systems can \emph{produce} accessible information, but whether they reliably \emph{interpret} patient wording across languages, phrasing variants, and care contexts \cite{recker-etal-2025-large,devries-etal-2025-enhancing}.
Despite this risk profile, much of the women's-health LLM literature still evaluates \emph{generation}: answer fluency, readability, expert preference, and user satisfaction \cite{adhikary-etal-2025-menstrual,recker-etal-2025-large}.
Domain adaptation can improve response quality on held-out questions, yet adaptation alone does not guarantee safe understanding.
Empirical work on perinatal patient information shows that LLM outputs vary with minor prompt phrasing, can diverge from national care guidelines, and may oversimplify complex decisions \cite{devries-etal-2025-enhancing,wang-etal-2024-prompt}.
Parallel lines of multilingual and cross-cultural NLP research further show that models often compress cultural variation, misalign with survey-grounded human judgments, and remain sensitive to surface form even when underlying intent is preserved \cite{mohammadi-etal-2025-large,nicholls-alperin-2025-cross}.
These findings motivate evaluation protocols that test \emph{interpretation}---condition inference, urgency calibration, clarification behavior, and cross-language consistency---rather than surface-level answer quality alone.

We introduce \textbf{HerHealthGPT-LU}, a multilingual benchmark for evaluating LLM \emph{understanding} of women's-health symptom communication across three categories: menstruation, polycystic ovary syndrome (PCOS)/hormonal concerns, and fertility.
Each of 90 English seed cases is expanded into six expression styles (canonical, clinical, layperson, indirect/culturally sensitive, ambiguous, and emotionally concerned), yielding 540 evaluation items whose clinical facts are held constant so that performance differences are attributable to linguistic form rather than content \cite{nicholls-alperin-2025-cross}.
Gold labels specify inferred condition, risk level, recommended action, and whether clarification is required, grounded where possible to public-health sources.
The benchmark is extended to Arabic and French through human-validated translation, enabling matched cross-lingual comparison.
Evaluation and adaptation are strictly separated: benchmark lineage is frozen, and fine-tuning draws only from a leakage-controlled silver corpus with dual exclusion by source provenance and paraphrase-family keys.

Our contributions are fourfold:
\begin{enumerate}
    \item A human-validated, multilingual benchmark of realistic women's-health symptom expressions spanning menstrual, PCOS, and fertility concerns, with style- and language-controlled item identities.
    \item An evaluation framework targeting misunderstanding, safety violations, clarification appropriateness, and cross-language/cross-style consistency---metrics aligned with clinical risk rather than lexical overlap alone.
    \item A systematic error taxonomy for analyzing how LLMs misinterpret women's-health language across models, languages, and registers.
    \item Evidence on whether multilingual LoRA fine-tuning on women's-health data reduces misunderstanding relative to general-purpose and domain-adapted baselines.
\end{enumerate}

The remainder of this paper reviews related work (\S\ref{sec:related}), describes the benchmark and experimental protocol (\S\ref{sec:methods}), presents results and error analysis (\S\ref{sec:results}), and discusses limitations and ethical considerations (\S\ref{sec:discussion}).

\section{Related Work}
\label{sec:related}

\subsection{Women's Health, Medical QA, and LLMs}

LLMs are increasingly positioned as digital companions for gynecological and obstetric patient education, offering anonymous, on-demand access to information on menstruation, contraception, pregnancy, and fertility \cite{recker-etal-2025-large}.
Editorial and empirical accounts emphasize both promise and hazard: models can simplify medical terminology and support health literacy, yet they lack individualized clinical assessment, may misalign with local guidelines, and can be misread as professional advice \cite{recker-etal-2025-large,devries-etal-2025-enhancing}.
De~Vries et al.\ report that ChatGPT can deliver relevant perinatal information across languages, but also document language-dependent accuracy and sensitivity to how questions are phrased \cite{devries-etal-2025-enhancing}.
These observations establish women's-health interaction as a setting where \emph{understanding} errors carry direct safety consequences, but they rarely provide benchmarked, item-level tests of interpretation across realistic patient registers.

Complementary work in obstetric patient education evaluates AI-generated responses on accuracy, comprehensiveness, and safety for specific conditions \cite{rotem-etal-2024-future}.
Such studies foreground expert-rated safety yet typically treat the patient query as given rather than systematically varying how the same clinical concern is expressed.
HerHealthGPT-LU addresses this gap by holding clinical content fixed while varying style and language, making linguistic understanding an explicit evaluation axis.

Prior PCOS-focused machine learning operates on structured clinical measurements---hormone panels, follicle counts, and tabular features---rather than free-text patient communication \cite{elmannai-etal-2023-polycystic}.
That line of work demonstrates the clinical importance of early hormonal-condition detection but does not evaluate whether language models interpret ambiguous symptom descriptions safely in conversation.
HerHealthGPT-LU complements these systems by targeting \emph{natural-language symptom communication}, where misunderstanding arises from wording rather than missing laboratory values.

\subsection{Domain Adaptation for Women's-Health LLMs}

MenstLLaMA is the closest prior \emph{system} work: a LoRA-adapted LLaMA-3-8B model trained on the MENST menstrual-health corpus and evaluated for educational answer quality in English \cite{adhikary-etal-2025-menstrual}.
The authors curate roughly 24K question--answer pairs, benchmark against general-purpose LLMs using automatic overlap metrics, expert ratings, and user satisfaction, and show gains in culturally sensitive menstrual-health education.
Crucially, however, MenstLLaMA optimizes and evaluates \emph{generation quality}---BLEU, BERTScore, expert preference---not whether a model correctly infers clinical intent, triage urgency, or clarification need from under-specified patient language.
MENST paraphrases also serve primarily as training-time augmentation without gold understanding labels tied to each surface form.
HerHealthGPT-LU uses MENST-derived material only under strict leakage control for adaptation, while treating style diversity as a \emph{controlled evaluation dimension} with gold condition, risk, action, and clarification labels.
The overlap in domain (menstrual health) is therefore superficial relative to task formulation: education-oriented generation versus multilingual understanding and safety evaluation across menstrual, PCOS, and fertility axes.

\subsection{Multilingual Clinical NLP and Cross-Cultural Evaluation}

The LUHME line of work studies whether LLMs preserve meaning and cultural variation beyond monolingual English settings.
Mohammadi et al.\ compare model-generated moral judgments with international survey data across dozens of countries, finding that LLMs often compress cross-cultural variance and achieve low alignment with empirical response patterns \cite{mohammadi-etal-2025-large}.
Their protocol---variance comparison, cluster alignment, significance testing on matched items---provides a template for rigorous cross-cultural evaluation, although their target is moral attitude probing rather than clinical symptom interpretation.
Nicholls and Alperin study cross-genre native language identification with open-source LLMs, showing that QLoRA fine-tuning transfers across registers when confounds such as topic-specific nouns are controlled \cite{nicholls-alperin-2025-cross}.
We adopt their leakage-aware adaptation philosophy and reporting discipline (multi-seed fine-tuning, per-class confusion analysis), but redirect the task from L1 classification to \emph{safe clinical understanding}: triage, clarification, and unsafe-advice detection.

HerHealthGPT-LU extends this multilingual evaluation tradition to women's-health \emph{patient language} in English, Arabic, and French.
Items are parallelized across languages with dual human validation to reduce translation-induced semantic drift---a known failure mode when symptom nuance and cultural phrasing alter clinical meaning \cite{mohammadi-etal-2025-large,recker-etal-2025-large}.

\subsection{Safety, Ambiguity, and Clarification}

Safe health communication requires models to recognize when a description is too vague for definitive guidance and to ask clarifying questions rather than over-commit \cite{recker-etal-2025-large,devries-etal-2025-enhancing}.
Prompt engineering studies likewise show that small wording changes can shift LLM behavior in clinical settings \cite{wang-etal-2024-prompt}.
Existing women's-health LLM evaluations seldom measure clarification appropriateness conditional on ambiguity, nor do they report severity errors such as under-triage of urgent presentations.
HerHealthGPT-LU labels items for required clarification and risk level, and scores models on misunderstanding rate, unsafe-response rate, and clarification behavior on ambiguous versus clear cases.

\subsection{Positioning of HerHealthGPT-LU}

Table~\ref{tab:positioning} summarizes how HerHealthGPT-LU differs from the closest prior lines of work.
Prior studies predominantly evaluate what models \emph{believe} \cite{mohammadi-etal-2025-large}, what they \emph{classify} \cite{nicholls-alperin-2025-cross,elmannai-etal-2023-polycystic}, or how well they \emph{generate} educational answers \cite{adhikary-etal-2025-menstrual,recker-etal-2025-large}.
We instead evaluate whether models \emph{understand patients safely}---interpreting multilingual, multi-register symptom communication with explicit tests of calibration, clarification, and cross-lingual consistency.
By separating a frozen gold benchmark from leakage-controlled fine-tuning data, HerHealthGPT-LU makes domain adaptation a hypothesis to be tested rather than a proxy for understanding.

\begin{table}[t]
\centering
\small
\begin{tabular}{lccc}
\hline
\textbf{Aspect} & \textbf{MenstLLaMA} & \textbf{LUHME'25 line} & \textbf{HerHealthGPT-LU} \\
\hline
Primary task & Answer generation & Attitudes / NLI & Symptom understanding \\
Languages & English & Mostly English & EN / AR / FR \\
Styles & Training aug. & Genre / prompt & Controlled eval axis \\
Safety metrics & Expert preference & Not clinical & Triage, clarify, unsafe \\
Benchmark gold & Held-out QA & Survey / NLI labels & Condition, risk, action \\
\hline
\end{tabular}
\caption{Positioning against closest related work. LUHME'25 line refers to \citet{mohammadi-etal-2025-large,nicholls-alperin-2025-cross}.}
\label{tab:positioning}
\end{table}

% --- Overleaf bibliography setup ---
% In acl_latex.tex (or main.tex), use:
%   \bibliography{custom,anthology-1,anthology-2}
% Add anthology shards via Overleaf: New File -> From External URL
%   https://aclanthology.org/anthology-1.bib
%   https://aclanthology.org/anthology-2.bib
% Keep custom.bib for non-ACL papers only; LUHME'25 proceedings resolve from anthology shards.
